\documentclass[12pt]{article}

\usepackage{amsmath, amssymb, amsthm}
\usepackage{enumerate}
\usepackage{hyperref}
\usepackage{geometry}
\usepackage{newtxtext}
\usepackage{newtxmath}
\geometry{margin=1in}

\title{Probability Theory - Lecture 2}
\author{}
\date{}

% ----------------------------------------------------------------
\begin{document}
\maketitle

One of the most basic principles in probability is:

\begin{center}
``The probability that at least one bad thing happens
is at most the sum of the probabilities of the bad things.''
\end{center}

This is called \textit{Boole's inequality}, or the \textit{union bound}.
It appears constantly in combinatorics, probability, computer science,
and additive combinatorics.

In this problem, we will derive it carefully from counting.

% ================================================================
\section*{Questions}

\begin{enumerate}[(a)]

\item Let $A$ and $B$ be finite sets.
Explain why
\[
|A\cup B|
=
|A|+|B|-|A\cap B|.
\]

\item Deduce that
\[
|A\cup B|
\leq
|A|+|B|.
\]

\item Let $A_1,A_2,\dots,A_n$ be finite sets.
Explain why
\[
\left|\bigcup_{i=1}^n A_i\right|
\leq
\sum_{i=1}^n |A_i|.
\]

(Hint: elements appearing in several sets get counted several times on the right-hand side.)

\item Suppose now that $S$ is a finite sample space with all outcomes equally likely.
If $E\subseteq S$ is an event, recall that
\[
\mathbb{P}(E)=\frac{|E|}{|S|}.
\]

Use part (c) to show that
\[
\mathbb{P}\left(\bigcup_{i=1}^n E_i\right)
\leq
\sum_{i=1}^n \mathbb{P}(E_i).
\]

\item Suppose a student takes $5$ exams.
For each exam, the probability of failing is at most $0.01$.

Show that the probability the student fails at least one exam is at most $0.05$.

\item A graph on $n$ vertices has edge set $E(X)$.
Suppose each edge independently survives with probability $p$.

For each edge $e$, let $E_e$ be the event that the edge survives.

Use Boole's inequality to show that the probability that at least one edge survives is at most
\[
|E(X)|p.
\]

\item Why can Boole's inequality be very wasteful when many events overlap heavily?

\item Give an example where equality holds in Boole's inequality.

\end{enumerate}

% ================================================================
\newpage
\section*{Solutions}

\begin{enumerate}[(a)]

\item When adding $|A|+|B|$, every element in $A\cap B$
gets counted twice:
once in $A$ and once in $B$.

Therefore we subtract $|A\cap B|$ once to correct this:
\[
|A\cup B|
=
|A|+|B|-|A\cap B|.
\]

\item Since
\[
|A\cap B|\geq 0,
\]
we have
\[
|A\cup B|
=
|A|+|B|-|A\cap B|
\leq
|A|+|B|.
\]

\item Every element belonging to at least one set $A_i$
is counted at least once in
\[
\sum_{i=1}^n |A_i|.
\]

If an element belongs to several sets,
it is counted several times on the right-hand side.

Hence the right-hand side can only be larger:
\[
\left|\bigcup_{i=1}^n A_i\right|
\leq
\sum_{i=1}^n |A_i|.
\]

\item By definition,
\[
\mathbb{P}(E_i)=\frac{|E_i|}{|S|}.
\]

Applying part (c),
\[
\left|\bigcup_{i=1}^n E_i\right|
\leq
\sum_{i=1}^n |E_i|.
\]

Dividing both sides by $|S|$ gives
\[
\mathbb{P}\left(\bigcup_{i=1}^n E_i\right)
=
\frac{\left|\bigcup E_i\right|}{|S|}
\leq
\frac{\sum |E_i|}{|S|}
=
\sum_{i=1}^n \mathbb{P}(E_i).
\]

This is Boole's inequality.

\item Let $E_i$ be the event that the student fails exam $i$.

Then
\[
\mathbb{P}(E_i)\leq 0.01
\]
for each $i$.

Hence by Boole's inequality,
\[
\mathbb{P}\left(\bigcup_{i=1}^5 E_i\right)
\leq
\sum_{i=1}^5 \mathbb{P}(E_i)
\leq
5(0.01)
=
0.05.
\]

\item Let
\[
F=\bigcup_{e\in E(X)} E_e
\]
be the event that at least one edge survives.

Applying Boole's inequality,
\[
\mathbb{P}(F)
\leq
\sum_{e\in E(X)} \mathbb{P}(E_e).
\]

Since each edge survives with probability $p$,
\[
\mathbb{P}(E_e)=p.
\]

Therefore
\[
\mathbb{P}(F)
\leq
\sum_{e\in E(X)} p
=
|E(X)|p.
\]

\item If many events overlap,
then the same outcomes get counted repeatedly in
\[
\sum_i \mathbb{P}(E_i).
\]

Thus the right-hand side may greatly overestimate the true probability of the union.

\item Equality holds when the events are pairwise disjoint.

For example, if
\[
E_1=\{1\},
\qquad
E_2=\{2\}
\]
inside the sample space
\[
S=\{1,2,3,4\},
\]
then
\[
\mathbb{P}(E_1\cup E_2)
=
\frac{2}{4}
=
\frac14+\frac14
=
\mathbb{P}(E_1)+\mathbb{P}(E_2).
\]
\qed

\end{enumerate}

\end{document}